% !PRINTED 130
\weberpage{130}
%<<CONT

and, denoting $\alpha+h$ by $\alpha'$,
\begin{equation}
 \alpha'=\alpha-\frac{a f(\alpha)}{2Q f'(\alpha)},
\tag{11}
\end{equation}
and, denoting $|f(\alpha')|$ by $a'$, obtain from (10)
\begin{equation}
 a'<a\left(1-\frac{a}{4Q}\right).
\tag{12}
\end{equation}

By the same procedure we derive from $\alpha',a'$ a second system $\alpha'',a''$, then from $\alpha'',a''$ a third system $\alpha''',a'''$, and so on; for the sequence of decreasing positive numbers $a,a',a'',\ldots$ we obtain the inequalities
\begin{equation}
\begin{aligned}
 a'&<a\left(1-\frac{a}{4Q}\right),\\
 a''&<a'\left(1-\frac{a'}{4Q}\right),\\
 a'''&<a''\left(1-\frac{a''}{4Q}\right),\\
 &\hspace{1.5em}\vdots
\end{aligned}
\tag{13}
\end{equation}
We now prove that this sequence approaches the limit zero. If this were not the case, a positive number $\omega$ could be given such that all $a^{(\nu)}$ remain above $\omega$; then the differences
\[
 \left(1-\frac{a}{4Q}\right),\quad
 \left(1-\frac{a'}{4Q}\right),\quad
 \left(1-\frac{a''}{4Q}\right)\cdots
\]
would all be smaller than the positive proper fraction
\[
 \Theta=1-\frac{\omega}{4Q},
\]
and from (13) it would follow that
\[
 a'<a\Theta,\quad
 a''<a'\Theta,\quad
 a'''<a''\Theta,\quad
 \ldots,\quad
 a^{(\nu)}<a^{(\nu-1)}\Theta.
\]
Whence, by multiplication,
\begin{equation}
 a^{(\nu)}<a\Theta^\nu.
\tag{14}
\end{equation}
This is a contradiction, since $\Theta^\nu$ becomes infinitely small as $\nu$ increases without bound. The result so obtained is expressed by
\begin{equation}
 \lim a^{(\nu)}=0.
\tag{15}
\end{equation}
